%% Description of the Jim-Tec / Arla butter cannon dataset.
%% Covers dimensions, combinations of faults/bearings/conditions, EDA findings,
%% supported tests and splits; the test rig is not described here.
%% Build (from eda/jim-tec/):   pdflatex jim-tec && pdflatex jim-tec

\documentclass[11pt,a4paper]{article}
\usepackage[T1]{fontenc}
\usepackage[utf8]{inputenc}
\usepackage[margin=2.4cm]{geometry}
\usepackage{amsmath}
\usepackage{booktabs}
\usepackage{graphicx}
\usepackage[hidelinks]{hyperref}

\title{Jim-Tec / Arla butter cannon}
\author{}
\date{\today}

\begin{document}
\maketitle

\section{Source}

\section{Dimensions of the time series}

\section{Combinations of faults, bearings and operating conditions}

\begin{table}[h]
\centering
\small
\begin{tabular}{ll}
\toprule
Property & Known \\
\midrule
Files         & CSV logs, 2026-04-14 to 2026-10-13 \\
Channels      & \texttt{Volt\_CH0} (sensor not documented), \texttt{RPM} \\
Sampling rate & 50\,kHz (logs with a header, 2026-10-12/13); unknown for older logs \\
Length        & 15\,min per headed log (ca.\ 45\,M samples) \\
Labels        & none \\
\bottomrule
\end{tabular}
\end{table}

\begin{itemize}
  \item No fault, bearing or operating-condition information beyond the logged RPM.
  \item RPM changes in steps of a few thousand samples (slower update rate).
  \item Source: \texttt{butter\_cannon\_eda.ipynb} (header checked by hand).
\end{itemize}

\section{Findings from EDA}

\section{Tests the data support}

\section{Splits}

\end{document}
